\pdfbookmark[1]{Abstract}{Abstract}
\chapter*{Abstract}

\textbf{The subject is the transform between two views of the same thing.} An object tracked through
a sequence of images, in any number of dimensions, is seen through a view that itself moves. What
the images record is the object's motion and the view's motion added together, and a tracker that
reads the sum as the object's motion follows the view instead of the object. This book collects the
published ways a view's motion is recovered from the images alone, in the order of the transforms
they recover: translation, rotation and scale, perspective, and flow over a curved surface that the
tracked objects travel on.

\textbf{What is quoted and what is new.} Every method in the first four chapters is published and
cited in the last, and is described here as its authors state it. Every one of them is computed in
floating point: a fast Fourier transform, a least squares fit, an iterative solver stopped at a
tolerance. The fifth chapter is the design this work adds, which carries the same transforms in
exact arithmetic: integer fields, number theoretic transforms whose counts never reach their
modulus, and rational rotations. Nothing in the fifth chapter has been measured here. It states what
exactness is predicted to change, so the predictions are on the page before any result is.

\textbf{Why the order matters.} A transform recovered first is removed before the next is looked
for, and each removal leaves a residual that the next transform explains. Plane plus parallax states
this for camera motion: register the dominant plane, and what remains is parallax plus the motion of
whatever moved on its own. The same stratification holds for any subject that drifts, turns and
deforms in front of a view that moves too, and it is the organising idea of the book.
